% =============================================
% File: chapters/introduction.tex
% MỞ ĐẦU — theo yêu cầu mục 3.7 HD/QT-PKHCN-NCKHSV
% =============================================

\chapter{MỞ ĐẦU}
\label{chap:modau}

% --------------------------------------------------
\section{Tổng quan tình hình nghiên cứu}
% --------------------------------------------------

Robot hai chân (bipedal robot) là một trong những hướng nghiên cứu trọng tâm của
lĩnh vực robot di động trong những thập kỷ gần đây. Các hệ thống như Boston Dynamics
Atlas, MIT Mini Cheetah hay Unitree H1 đã chứng minh tiềm năng ứng dụng của robot
bipedal trong môi trường phức tạp, nơi robot bánh xe truyền thống không thể hoạt động
hiệu quả. Tuy nhiên, chi phí nghiên cứu và phát triển những hệ thống này còn rất cao,
gây khó khăn cho các nhóm nghiên cứu quy mô nhỏ và sinh viên.

Ở phân khúc thiết kế đơn giản hóa hơn, mô hình \textit{bipedal wheel robot} —
robot hai chân có bánh xe tại điểm cuối mỗi chân — đang được quan tâm vì cân bằng
được giữa tính linh hoạt của chân và hiệu quả di chuyển của bánh xe. Điển hình là các
nghiên cứu về robot Wheeled-Bipedal tại ETH Zurich và các biến thể sử dụng cơ cấu
song song để tăng cứng vững và giảm số bậc tự do cần điều khiển.

Trong lĩnh vực điều khiển, bộ điều khiển PID (Proportional–Integral–Derivative)
vẫn là lựa chọn phổ biến nhờ tính đơn giản và hiệu quả. Song, PID truyền thống có hệ
số cố định, dẫn đến suy giảm hiệu năng khi hệ thống hoạt động ở nhiều dải tải và vận
tốc khác nhau. Sự kết hợp giữa PID và Mạng Nơron Nhân Tạo (ANN) để dự đoán và
hiệu chỉnh tham số điều khiển thích nghi đang là xu hướng được nghiên cứu rộng rãi,
cho thấy kết quả cải thiện đáng kể về độ vọt lố và thời gian xác lập.

Tại Việt Nam, các nghiên cứu về điều khiển thích nghi ứng dụng mạng nơron cho
hệ truyền động cơ DC còn khá hạn chế ở quy mô sinh viên. Việc xây dựng một hệ
thống hoàn chỉnh từ nhận dạng mô hình động cơ, thiết kế bộ PI, huấn luyện ANN đến
triển khai thực tế trên phần cứng nhúng là bài toán còn nhiều khoảng trống để nghiên cứu
và đóng góp.

% --------------------------------------------------
\section{Tính cấp thiết của đề tài}
% --------------------------------------------------

Trong bối cảnh Cách mạng công nghiệp 4.0, các lĩnh vực tự động hóa, robot thông minh
và trí tuệ nhân tạo đang phát triển mạnh mẽ, tạo ra nhiều hướng nghiên cứu mới trong
thiết kế và điều khiển các hệ thống robot di động. Trong đó, robot tự cân bằng là một
hướng nghiên cứu quan trọng vì có khả năng duy trì ổn định trong quá trình di chuyển,
đồng thời làm nền tảng cho các hệ robot tự hành, robot dịch vụ và phương tiện di chuyển
thông minh.

Hiện nay, phần lớn các mô hình robot tự cân bằng vẫn tập trung vào dạng hai bánh truyền
thống và sử dụng các thuật toán điều khiển như PID, LQR hoặc Fuzzy Logic. Các phương pháp
này có ưu điểm là đơn giản và dễ triển khai, tuy nhiên còn hạn chế khi hệ thống có nhiều
bậc tự do, đặc tính phi tuyến mạnh hoặc chịu tác động của môi trường thay đổi. Vì vậy,
việc nghiên cứu xe robot tự cân bằng đa bậc tự do kết hợp với trí tuệ nhân tạo và học
tăng cường có ý nghĩa quan trọng, giúp robot có khả năng tự học, thích nghi và nâng cao
độ ổn định trong quá trình vận hành.

Đề tài ``Xe robot tự cân bằng đa bậc tự do'' được thực hiện nhằm xây dựng nền tảng nghiên
cứu cho hệ robot di động thông minh, kết hợp giữa mô hình cơ khí, cảm biến, điều khiển
nhúng, mô phỏng vật lý và thuật toán AI. Kết quả của đề tài không chỉ có ý nghĩa trong
nghiên cứu học thuật mà còn có tiềm năng ứng dụng trong robot dịch vụ, robot giao hàng,
xe tự hành mini và các mô hình đào tạo về điều khiển -- trí tuệ nhân tạo. Đồng thời, đề
tài giúp sinh viên rèn luyện tư duy thiết kế hệ thống, lập trình điều khiển và tối ưu
hóa robot trong thời kỳ chuyển đổi số.

% --------------------------------------------------
\section{Lý do chọn đề tài}
% --------------------------------------------------

Xuất phát từ nhu cầu nghiên cứu và phát triển các hệ robot di động thông minh có khả
năng tự cân bằng, thích nghi và hoạt động ổn định trong nhiều điều kiện khác nhau,
nhóm nghiên cứu chọn đề tài:

\begin{center}
    \textbf{``Xe robot tự cân bằng đa bậc tự do''}
\end{center}

với các lý do cụ thể sau:

\begin{enumerate}[leftmargin=1.5cm, label=\arabic*.]
    \item \textbf{Tính thực tiễn cao:} Robot tự cân bằng là nền tảng quan trọng cho
    các hệ thống robot di động, xe tự hành mini, robot dịch vụ và robot hỗ trợ con người.
    Việc nghiên cứu mô hình xe robot tự cân bằng đa bậc tự do giúp tạo cơ sở để phát triển
    các hệ robot linh hoạt, có khả năng di chuyển ổn định và thích nghi tốt hơn.

    \item \textbf{Tính học thuật:} Đề tài kết hợp nhiều kiến thức quan trọng như mô hình
    động học, động lực học, cảm biến quán tính, điều khiển cân bằng, điều khiển động cơ,
    trí tuệ nhân tạo và học tăng cường. Qua đó, đề tài giúp xây dựng nền tảng nghiên cứu
    toàn diện cho hệ robot phi tuyến và nhiều bậc tự do.

    \item \textbf{Tính mới và khả năng mở rộng:} Khác với các mô hình robot tự cân bằng
    hai bánh truyền thống, robot tự cân bằng đa bậc tự do có khả năng mở rộng về cơ cấu,
    chuyển động và chiến lược điều khiển. Việc kết hợp mô phỏng vật lý với học tăng cường
    giúp robot có khả năng tự học cách cân bằng và di chuyển trong môi trường mô phỏng
    trước khi triển khai trên phần cứng thực tế.

    \item \textbf{Đóng góp vào nghiên cứu trong nước:} Đề tài góp phần xây dựng một hướng
    nghiên cứu phù hợp với điều kiện sinh viên Việt Nam, kết hợp giữa thiết kế cơ khí,
    điều khiển nhúng, mô phỏng và trí tuệ nhân tạo. Kết quả nghiên cứu có thể làm nền tảng
    cho các đề tài tiếp theo về robot di động thông minh, robot dịch vụ và hệ thống tự hành.
\end{enumerate}


% --------------------------------------------------
\section{Mục tiêu đề tài}
% --------------------------------------------------

Đề tài hướng đến các mục tiêu cụ thể sau:

\begin{enumerate}[leftmargin=1.5cm, label=\arabic*.]
    \item Nhận dạng mô hình hàm truyền của động cơ DC NF5475 thông qua thu thập
    dữ liệu thực nghiệm và công cụ \texttt{tfest} trong MATLAB.

    \item Thiết kế bộ điều khiển PI vận tốc với tham số tối ưu sử dụng giải thuật di
    truyền (GA), đảm bảo độ vọt lố dưới 5\% và thời gian xác lập dưới 0{,}5s.

    \item Xây dựng và huấn luyện mạng ANN có khả năng dự đoán tham số $K_p$, $K_i$
    thích nghi theo trạng thái sai số của hệ thống, giảm thiểu sai số bám setpoint khi
    vận tốc đặt thay đổi.

    \item Triển khai toàn bộ hệ thống điều khiển PI--ANN trên vi điều khiển ESP32
    WROOM, kiểm chứng hiệu năng thực nghiệm.

    \item Phân tích động học (thuận và nghịch) của cơ cấu 5-bar parallel, xác định
    không gian làm việc và Jacobian, làm nền tảng cho điều khiển quỹ đạo điểm cuối
    chân robot.
    
    \item Triển khai thuật toán học tăng cường (Reinforcement Learning) nhằm giúp mô hình học được khả năng tự cân bằng và di chuyển trong quá trình vận hành.
\end{enumerate}


% --------------------------------------------------
\section{Phương pháp nghiên cứu}
% --------------------------------------------------

Đề tài sử dụng kết hợp các phương pháp sau:

\begin{itemize}[leftmargin=1.5cm]
    \item \textbf{Phương pháp lý thuyết:} Nghiên cứu tài liệu về lý thuyết điều khiển
    PID, mạng nơron nhân tạo (ANN), động học cơ cấu song song và nhận dạng hệ thống.

    \item \textbf{Phương pháp thực nghiệm:} Thu thập dữ liệu đặc tính động cơ trực
    tiếp trên phần cứng; kiểm chứng bộ điều khiển qua các kịch bản thực nghiệm với
    các mức vận tốc đặt khác nhau (250 RPM đến 2250 RPM).

    \item \textbf{Phương pháp mô phỏng:} Sử dụng MATLAB/Simulink để nhận dạng hàm
    truyền, thiết kế và kiểm tra bộ điều khiển PI trước khi nạp lên phần cứng; sử dụng
    Python để huấn luyện và đánh giá mạng ANN; sử dụng Isaac Sim/Isaac Lab để mô phỏng
    môi trường vật lý và huấn luyện học tăng cường trước khi triển khai trên phần cứng.

    \item \textbf{Phương pháp tối ưu hóa:} Áp dụng giải thuật di truyền (Genetic
    Algorithm) để tìm bộ tham số $K_p$, $K_i$ tối ưu trong không gian ràng buộc. Đồng thời,
    tận dụng khả năng tính toán của GPU để huấn luyện học tăng cường song song nhiều môi trường mô phỏng,
    nhằm rút ngắn thời gian huấn luyện và đánh giá thuật toán.
\end{itemize}


% --------------------------------------------------
\section{Đối tượng và phạm vi nghiên cứu}
% --------------------------------------------------

\textbf{Đối tượng nghiên cứu:}
\begin{itemize}[leftmargin=1.5cm]
    \item Động cơ DC Brushed NF5475 (12--24V, 33W, encoder 200 PPR).
    \item Vi điều khiển ESP32 WROOM-32 lập trình bằng Arduino IDE.
    \item Mạng Nơron Nhân Tạo (ANN) xây dựng bằng Python không dùng framework.
    \item Cơ cấu 5-bar parallel với thông số: $A_1A_2 = 105$~mm,
    $L_{act} = 100$~mm, $L_{pas} = 180$~mm.
    \item Thuật toán học tăng cường sử dụng thuật toán Proximal Policy Optimization (PPO),
    được mô phỏng và huấn luyện trên phần mềm mô phỏng Isaac Sim kết hợp với nền tảng Isaac Lab.
\end{itemize}

\textbf{Phạm vi nghiên cứu:}
\begin{itemize}[leftmargin=1.5cm]
    \item Điều khiển vận tốc một động cơ DC trong vòng lặp kín, dải setpoint từ
    250 đến 2250 RPM.
    \item ANN dự đoán hai tham số $K_p$ và $K_i$; đầu vào gồm sai số vận tốc
    ($e$) và đạo hàm sai số ($\Delta e$).
    \item Phân tích động học phẳng 2D của cơ cấu 5-bar parallel (chưa bao gồm
    động lực học và điều khiển toàn hệ robot).
    \item Thực nghiệm được thực hiện trong môi trường phòng thí nghiệm, tải cố
    định, chưa xét đến nhiễu từ ngoài.
\end{itemize}
